\section{Locality and maxima of bounded potentials}
\label{sec:hessian-bounded-locality}

Where one bounded admissible potential is strictly larger than another,
the maximum has the same positive Hessian current as the larger
potential. These strict comparison sets need not be open.
Quasicontinuity gives cutoffs whose products with the approximating
currents converge, allowing smooth locality to pass to these sets.

\Needspace{15\baselineskip}
\begin{proposition}\label{hessian:bounded-strict-locality}\index{Hessian measure!locality}
Let $u,v$ be bounded admissible, put $w=\max\{u,v\}$, and fix
$1\leq p\leq m$. For any bounded admissible $t_2,\ldots,t_p$,
\[
 \mathbf1_{\{u>v\}}Q_p(w,t_2,\ldots,t_p)
       =\mathbf1_{\{u>v\}}Q_p(u,t_2,\ldots,t_p).
\]
The equality is coefficientwise as currents with measure
coefficients. The corresponding equality on $\{v>u\}$ follows
by exchanging $u,v$.
\end{proposition}
\begin{proof}
Choose smooth uniformly bounded decreasing approximations of all
the potentials by \cref{hessian:bounded-decreasing-approximation}.
Construct $w_j$ by \cref{hessian:regularized-maximum}. Work on
one compact coordinate test set. Fix $\delta>0$ and a smooth
$\chi_\delta:\mathbb R\to[0,1]$ that vanishes for
$s\leq\delta/2$ and equals one for $s\geq\delta$.
For $j$ large enough that $\epsilon_j<\delta/2$, ordinary smooth
locality gives
\[
 \chi_\delta(u_j-v_j)Q_p(w_j,t_2^j,\ldots,t_p^j)
       =\chi_\delta(u_j-v_j)Q_p(u_j,t_2^j,\ldots,t_p^j).
\]
Indeed, wherever the weight is nonzero, $w_j=u_j$ on an open
neighborhood. The equality elsewhere follows from its zero weight.

Both sequences of positive currents converge by
\cref{hessian:joint-decreasing-products}; every coefficient variation
is uniformly dominated by $C c_m^*$ by
\cref{hessian:positive-capacity-domination}. The potentials and their
limits are quasicontinuous by \cref{hessian:rough-coefficient-capacity}.
Apply \cref{hessian:weighted-measure-convergence} with $s=m$ and
$F(x,y)=\chi_\delta(x-y)$ to the real and imaginary parts of
each coefficient. Passing to the limit gives the same identity
with $u,v,w,t_2,\ldots,t_p$.
Choose the cutoffs so that $\chi_\delta(s)\to\mathbf1_{\{s>0\}}$
pointwise as $\delta\downarrow0$. Dominated convergence against
the fixed finite coefficient variations on the test set gives
the asserted restriction. A partition of unity proves it globally.
\end{proof}

On the contact set, a positive constant shift gives a strict comparison
before the limit is taken. Decreasing the shift then recovers that
set by dominated convergence for a fixed measure.

\Needspace{16\baselineskip}
\begin{theorem}\label{hessian:bounded-maximum-measure}\index{Hessian measure!maximum inequality}
Let $u,v$ be bounded admissible. Fix $1\leq r\leq m$ and bounded admissible
$z_1,\ldots,z_{m-r}$. Define the positive Radon measure
\[
 \mu_r(t)=Q_m(\underbrace{t,\ldots,t}_{r},z_1,\ldots,z_{m-r}).
\]
For $w=\max\{u,v\}$,
\[
 \mu_r(w)\geq\mathbf1_{\{u\geq v\}}\mu_r(u)
                  +\mathbf1_{\{u<v\}}\mu_r(v).
\]
The inequality with $u,v$ exchanged also holds.
\end{theorem}
\begin{proof}
Symmetry and \cref{hessian:bounded-strict-locality}, applied to the $r$ slots, give
$\mu_r(w)=\mu_r(v)$ on $\{u<v\}$ and $\mu_r(w)=\mu_r(u)$ on $\{u>v\}$.
To include the contact set, put $w_\epsilon=\max\{u+\epsilon,v\}$.
Strict locality and invariance under addition of constants give
\[
 \mu_r(w_\epsilon)\geq
                   \mathbf1_{\{u+\epsilon>v\}}\mu_r(u).
\]
The invariance follows from the smooth products and their joint
decreasing limits. As $\epsilon\downarrow0$,
$w_\epsilon\downarrow w$, so the left side converges weakly to
$\mu_r(w)$. Against any nonnegative continuous compactly supported
test function, the right side converges by dominated convergence
for the fixed measure $\mu_r(u)$ to
$\mathbf1_{\{u\geq v\}}\mu_r(u)$. This proves the measure inequality
on the contact set as well as on $\{u>v\}$. On its disjoint
complement $\{u<v\}$, the already proved equality supplies
$\mu_r(v)$. Restricting the positive measures to these two Borel sets
and adding gives the assertion. Exchanging $u,v$ gives its other
form.
\end{proof}
